\documentclass{article}
%DIF LATEXDIFF DIFFERENCE FILE
%DIF DEL cover_letter_old_temp.tex   Fri Mar  6 18:51:46 2026
%DIF ADD cover_letter_new_temp.tex   Fri Mar  6 18:51:46 2026

\usepackage{siunitx}
\usepackage{hyperref}
\usepackage[english]{datetime2}
%DIF PREAMBLE EXTENSION ADDED BY LATEXDIFF
%DIF UNDERLINE PREAMBLE %DIF PREAMBLE
\RequirePackage[normalem]{ulem} %DIF PREAMBLE
\RequirePackage{color}\definecolor{RED}{rgb}{1,0,0}\definecolor{BLUE}{rgb}{0,0,1} %DIF PREAMBLE
\providecommand{\DIFaddtex}[1]{{\protect\color{blue}\uwave{#1}}} %DIF PREAMBLE
\providecommand{\DIFdeltex}[1]{{\protect\color{red}\sout{#1}}} %DIF PREAMBLE
%DIF SAFE PREAMBLE %DIF PREAMBLE
\providecommand{\DIFaddbegin}{} %DIF PREAMBLE
\providecommand{\DIFaddend}{} %DIF PREAMBLE
\providecommand{\DIFdelbegin}{} %DIF PREAMBLE
\providecommand{\DIFdelend}{} %DIF PREAMBLE
\providecommand{\DIFmodbegin}{} %DIF PREAMBLE
\providecommand{\DIFmodend}{} %DIF PREAMBLE
%DIF FLOATSAFE PREAMBLE %DIF PREAMBLE
\providecommand{\DIFaddFL}[1]{\DIFadd{#1}} %DIF PREAMBLE
\providecommand{\DIFdelFL}[1]{\DIFdel{#1}} %DIF PREAMBLE
\providecommand{\DIFaddbeginFL}{} %DIF PREAMBLE
\providecommand{\DIFaddendFL}{} %DIF PREAMBLE
\providecommand{\DIFdelbeginFL}{} %DIF PREAMBLE
\providecommand{\DIFdelendFL}{} %DIF PREAMBLE
%DIF HYPERREF PREAMBLE %DIF PREAMBLE
\providecommand{\DIFadd}[1]{\texorpdfstring{\DIFaddtex{#1}}{#1}} %DIF PREAMBLE
\providecommand{\DIFdel}[1]{\texorpdfstring{\DIFdeltex{#1}}{}} %DIF PREAMBLE
%DIF COLORLISTINGS PREAMBLE %DIF PREAMBLE
\RequirePackage{listings} %DIF PREAMBLE
\RequirePackage{color} %DIF PREAMBLE
\lstdefinelanguage{DIFcode}{ %DIF PREAMBLE
%DIF DIFCODE_UNDERLINE %DIF PREAMBLE
  moredelim=[il][\color{red}\sout]{\%DIF\ <\ }, %DIF PREAMBLE
  moredelim=[il][\color{blue}\uwave]{\%DIF\ >\ } %DIF PREAMBLE
} %DIF PREAMBLE
\lstdefinestyle{DIFverbatimstyle}{ %DIF PREAMBLE
	language=DIFcode, %DIF PREAMBLE
	basicstyle=\ttfamily, %DIF PREAMBLE
	columns=fullflexible, %DIF PREAMBLE
	keepspaces=true %DIF PREAMBLE
} %DIF PREAMBLE
\lstnewenvironment{DIFverbatim}{\lstset{style=DIFverbatimstyle}}{} %DIF PREAMBLE
\lstnewenvironment{DIFverbatim*}{\lstset{style=DIFverbatimstyle,showspaces=true}}{} %DIF PREAMBLE
\lstset{extendedchars=\true,inputencoding=utf8}

%DIF END PREAMBLE EXTENSION ADDED BY LATEXDIFF

\begin{document}

\today

Dear Editors of Nano Letters:

Please find enclosed our manuscript titled ``Water Does Not Slip on Hydrophobic Surfaces: Insights from Slip Length Mapping" by Haruya Ishida, Hideaki Teshima, Koji Takahashi, Vishwanath Ganesan, and Nenad Miljkovic, for consideration as a Letter in Nano Letters.

A persistent problem in interfacial physics is that experimentally reported slip lengths at solid-liquid interfaces span orders of magnitude even for nominally similar systems. This wide variability contrasts sharply with molecular dynamics simulations, which predict well-defined and nearly universal trends, and has obscured the interpretation of what measured slip lengths actually represent. A central difficulty is the lack of experimental access to the intrinsic slip boundary condition, disentangled from nanoscale surface features that can masquerade as slip.

Here, we address this long-standing challenge by developing a frequency-modulation AFM approach that boosts hydrodynamic (slip) sensitivity by $159\times$ and enables co-registered nanoscale mapping of slip length and surface topography. This simultaneous mapping is crucial: it allows us to directly access intrinsic \DIFdelbegin \DIFdel{(``true'') }\DIFdelend slip on genuinely flat regions while identifying and excluding apparent slip arising from unavoidable nanoscale heterogeneities. In a single measurement, we can determine where slip is intrinsic, where it is an artifact, and why.

Our key findings are summarized as follows:
\begin{enumerate}
\item Most ``hydrophobic slip" is not intrinsic slip. Using nanoscopically flat substrates spanning intrinsic contact angles from \SI{3}{\degree} to \SI{111}{\degree}, we find that the \DIFdelbegin \DIFdel{true }\DIFdelend \DIFaddbegin \DIFadd{intrinsic }\DIFaddend slip length on \DIFdelbegin \DIFdel{them }\DIFdelend \DIFaddbegin \DIFadd{these surfaces }\DIFaddend is essentially zero (sub-nanometer; even Teflon: $0.8 \pm  \SI{0.8}{\nano\metre}$). Our dataset quantitatively follows the molecular-dynamics-predicted wettability scaling, thereby reconciling a decades-old mismatch between simulations and experiments.

\item Apparent slip can be directly identified---and removed. Our mapping shows that nanoscale geometry and nanobubbles generate pronounced local changes in the apparent \DIFdelbegin \DIFdel{slips}\DIFdelend \DIFaddbegin \DIFadd{slip}\DIFaddend . This is the key mechanism behind the notorious irreproducibility in the literature: the ``slip length'' many techniques report is often a convolution of boundary condition and nanoscale geometry.

\item Graphite \DIFdelbegin \DIFdel{is }\DIFdelend \DIFaddbegin \DIFadd{appears to be }\DIFaddend a genuine exception, but only under specific conditions. On graphite in deionized water, we measure a large slip length of $43.2 \pm \SI{5.8}{\nano\metre}$---yet this slippage \DIFdelbegin \DIFdel{collapses to a }\DIFdelend \DIFaddbegin \DIFadd{decreases to a near }\DIFaddend no-slip condition in electrolyte solutions. This striking on/off behavior highlights the decisive influence of interfacial charge/ion adsorption in pinning water at otherwise ultra-smooth, chemically homogeneous carbon interfaces.
\end{enumerate}

More broadly, our co-registered FM-AFM approach turns the slip boundary condition from a fitted parameter into a directly imaged, quantitative field variable. This platform enables systematic tests of how ions/charge, soft or porous interfaces, and nanoscale gas phases regulate interfacial transport, providing reproducible benchmarks for theory and guiding predictive nanofluidic design. We believe this manuscript will be of strong interest to the readership of Nano Letters, given its emphasis on nanoscale physics of fluids, surfaces and interfaces, and transport phenomena. Given the ongoing debate and reproducibility issues surrounding interfacial slip, we believe these results are time-sensitive and warrant rapid communication.

We confirm that this manuscript has not been published elsewhere and is not under consideration in whole or in part by another journal. The authors declare no competing interests.

We would like to recommend the following potential referees: 
\begin{itemize}
\item Prof. Lydﾃｩric Bocquet\\
(ﾃ営ole Normale Supﾃｩrieure, France; lyderic.bocquet@ens.fr) 
\item Prof. Timothﾃｩe Mouterde\\
(The University of Tokyo, Japan; mouterde@g.ecc.u-tokyo.ac.jp) 
\item Prof. Laurent Joly\\
(Universitﾃｩ Lyon 1, France; laurent.joly@univ-lyon1.fr)
\item Prof. Enrique Wagemann\\
(Universidad de Concepciﾃｳn, Chile; ewagemann@udec.cl)
\item Prof. Sushanta Mitra\\
(University of Waterloo, Canada; skmitra@uwaterloo.ca)
\item Prof. Daniel Orejon\\
(The University of Edinburgh, UK; D.Orejon@ed.ac.uk)
\item Prof. Daniel J. Preston\\
(Rice University, USA, djp@rice.edu)
\item Prof. Shalabh Maroo\\
(Syracuse University, USA, scmaroo@syr.edu)
\item Prof. Dion Antao\\
(Texas A\&M University, USA, dantao@tamu.edu)
\end{itemize}

Yours sincerely,\\
Hideaki Teshima (Corresponding author)\\
Associate Professor\\
Department of Aeronautics and Astronautics\\
Kyushu University\\
Nishi-Ku, Motooka 744 Fukuoka 819-0395, Japan\\
Tel: +81-92-802-3050\\
E-mail: hteshima05@aero.kyushu-u.ac.jp

\end{document}
